\section{System Design}
\label{sec:design}

\subsection{Roles}
\begin{description}[leftmargin=2.2cm,style=nextline,itemsep=2pt]
  \item[Organizer] creates a circle and sets its terms; may cancel it while it is still open. Never
        touches funds.
  \item[Member] joins with a deposit, takes part in the payout draw, pays every round, receives the pot once.
  \item[Admin] the platform operator (the deploying account). May pause joining and paying in an
        emergency, for a bounded time (Section~\ref{sec:admin}); cannot move money.
  \item[Anyone] may finalize a draw, settle a round after its deadline, or cancel a circle that never
        filled. The keeper script does this automatically, but it is a convenience, not a trust point.
\end{description}

\subsection{Workflows}
A circle runs eight workflows (the brief asks for at least three); every one changes real state.

\begin{table}[H]
\centering\small
\caption{User workflows and the functions behind them}
\label{tab:workflows}
\begin{tabularx}{\textwidth}{@{}P{0.55cm}P{2.3cm}P{1.9cm}LL@{}}
\toprule
\# & Workflow & Actor & On chain & Off chain \\
\midrule
W1 & Create circle & Organizer & \code{CircleFactory.\allowbreak createCircle}: clone, initialize and register a \code{Circle} with amount, size, durations and terms CID & Terms uploaded to IPFS first; the browser checks the CID \\
W2 & Join & Member & approve the vault, \code{Circle.join}: lock the slot 1 deposit; reputation checked; the last join starts the draw & Display name, signed by the wallet, stored in the indexer \\
W3 & Draw payout order & All members & \code{PayoutOrder.commit}, \code{reveal}; \code{Circle.finalizeDraw} sets the order, removes non revealers, refunds excess deposits & Reminders before commit and reveal deadlines \\
W4 & Contribute & Member & approve the circle, \code{Circle.contribute}: paid bit set & Due date reminders \\
W5 & Settle round & Anyone (keeper) & \code{Circle.settleRound}: slash unpaid members, record defaults, credit the pot & Keeper calls it after the deadline \\
W6 & Claim payout & Recipient & \code{Circle.claimPayout}: pull payment of every credit & \enquote{Your pot is ready} notification \\
W7 & Complete & Automatic & last \code{settleRound}: leftover debts, release deposits, record completions & History from the indexer \\
W8 & Cancel & Organizer, or anyone after the join deadline & \code{Circle.cancel}: everyone refunded & Notification to joined members \\
\bottomrule
\end{tabularx}
\end{table}

\subsection{Lifecycle}
Figure~\ref{fig:lifecycle} shows the five states and the three steps of every round. Any failure
before the rounds start ends in a full refund; a missed payment during the rounds is covered by
slashing, not by stopping the circle. \code{BaseCircleLifecycle} enforces the transitions: every
state changing function has an \code{inState} guard and \code{\_transition} rejects any edge not in the diagram.

\begin{figure}[H]
\centering
\begin{tikzpicture}[
  st/.style={draw=boxline, rounded corners=3pt, align=center, minimum width=3.0cm, minimum height=1.15cm, font=\small, fill=white},
  ar/.style={-{Stealth[length=2.2mm]}, thick, draw=black!75},
  lb/.style={font=\scriptsize, align=center}]
\node[st, fill=tintblue] (open) at (0,0) {\textbf{Open}\\members join and\\lock a deposit};
\node[st, fill=tintblue] (draw) at (4.7,0) {\textbf{Order draw}\\commit, reveal;\\excess refunded};
\node[st] (contrib) at (10.0,0.2) {\textbf{Contributing}\\pay before deadline\\plus grace};
\node[st] (settle) at (10.0,-1.7) {\textbf{Settle round}\\unpaid members\\slashed and logged};
\node[st] (payout) at (10.0,-3.6) {\textbf{Payout claimable}\\recipient pulls\\the pot};
\node[st, fill=tintgreen, line width=1.1pt] (done) at (10.0,-5.9) {\textbf{Completed}\\deposits released,\\reputation updated};
\node[st, fill=tintrose] (cancel) at (0,-3.6) {\textbf{Cancelled}\\everyone refunded\\in full};
\begin{scope}[on background layer]
  \node[draw=boxline, dashed, rounded corners=5pt, fill=tintgreen!70, inner sep=7pt, fit=(contrib)(settle)(payout), label={[font=\small\bfseries]above:Active: rounds 1 to N}] (rounds) {};
\end{scope}
\draw[ar] (open) -- node[lb, above]{last member\\joins} (draw);
\draw[ar] (draw) -- node[lb, above]{2 or more\\revealed} (contrib.west |- draw);
\draw[ar] (contrib) -- (settle);
\draw[ar] (settle) -- (payout);
\draw[ar] (payout.east) -- ++(0.75,0) |- node[lb, right, pos=0.25]{next\\round} (contrib.east);
\draw[ar] (payout) -- node[lb, right]{after round N} (done);
\draw[ar] (open) -- node[lb, left]{organizer, or\\join deadline\\passed unfilled} (cancel);
\draw[ar] (draw.south) |- node[lb, above, pos=0.75]{fewer than 2 revealed} (cancel.east);
\end{tikzpicture}
\caption{Circle lifecycle: five states, three steps per round}
\label{fig:lifecycle}
\end{figure}

\subsection{Collateral rule}
\label{sec:collateral}
A member at payout position $k$ in a circle of $N$ participants still owes $N-k$ contributions after
receiving the pot. The vault keeps exactly that much locked:
\begin{equation}
  \text{requiredCollateral}(c, N, k) = c \times (N - k)
\end{equation}
where $c$ is the contribution. Coverage is 100\%: our planning review found that the originally
planned 50\% would leave the group short by half of every missed payment, which contradicts
\enquote{others are made whole}. Early slots lock more and the last slot locks nothing, which keeps
the credit benefit that flat collateral loses. To stay simple, everyone locks the slot 1 amount
$c(N-1)$ on joining and the vault refunds the excess once the draw fixes each position. For 4
members paying 1,000 mUSD the deposits after the draw are 3,000, 2,000, 1,000 and 0 mUSD.

When a round is settled, for every member who did not pay:
\begin{itemize}
  \item a member who already received the pot is slashed one contribution (always fully covered);
  \item a member not yet paid is slashed what their deposit holds; any uncovered part becomes a debt
        (\code{debtOf}) owed to that round's recipient (\code{deficitOf});
  \item a member who misses their own payout round simply receives a smaller pot.
\end{itemize}
When the recipient is paid, any debt they owe is repaid out of their pot first and credited to whoever
was shortchanged, and their deposit is refilled from the pot back to 100\% of what they still owe.
Section~\ref{sec:security} proves this always suffices (attack test A8).

\subsection{Architecture}
Figure~\ref{fig:architecture} shows the components and the trust boundary. Transactions go straight
from the member's wallet to the contracts; the off chain services are never in the path of a money
moving transaction and only read from the chain (the keeper only calls functions anyone may call).

\begin{figure}[H]
\centering
\begin{tikzpicture}[
  box/.style={draw=boxline, rounded corners=2pt, align=center, font=\small, minimum height=0.95cm, inner sep=3pt},
  pp/.style={box, fill=people, minimum width=2.5cm},
  oc/.style={box, fill=onchain, minimum width=2.3cm, font=\footnotesize},
  of/.style={box, fill=offchain, minimum width=3.1cm},
  ar/.style={-{Stealth[length=2mm]}, draw=black!75},
  lb/.style={font=\scriptsize, align=center, fill=white, inner sep=1pt}]
% people
\node[pp] (org) at (-3.6,6.9) {Organizer};
\node[pp] (mem) at (0.3,6.9) {Members};
\node[pp] (adm) at (4.2,6.9) {Admin (pause only)};
% browser
\node[box, fill=white, minimum width=13.6cm, minimum height=1.1cm] (ui) at (0.3,5.3) {\textbf{Browser app} (React, Vite, ethers v6): wallet, role based views, transaction tray, error decoder,\\terms re hash check, chain fallback when the indexer is down};
% off chain
\node[of] (ipfs) at (-4.3,2.9) {IPFS terms store\\(CIDv1, sha256)};
\node[of] (idx) at (-4.3,1.5) {Indexer + HTTP API\\(SQLite, rebuildable)};
\node[of] (noti) at (-4.3,0.1) {Notifier\\(in app reminders)};
\node[of] (keep) at (-4.3,-1.3) {Keeper\\(settle, finalize, cancel)};
% on chain
\node[oc, minimum width=3.1cm, font=\small] (fac) at (2.6,2.9) {\code{CircleFactory}};
\node[oc, minimum width=2.4cm] (tok) at (6.4,2.9) {\code{MockStablecoin}};
\node[oc, minimum width=2.8cm, font=\small] (circ) at (2.6,1.5) {\code{Circle} clones\\(one per committee)};
\node[oc, minimum width=2.4cm] (impl) at (6.4,1.5) {\code{Circle}\\implementation};
\node[oc] (vault) at (0.6,-0.4) {\code{Collateral}\\\code{Vault}};
\node[oc] (order) at (3.1,-0.4) {\code{PayoutOrder}};
\node[oc] (rep) at (5.6,-0.4) {\code{Reputation}\\\code{Registry}};
\begin{scope}[on background layer]
  \node[draw=boxline, rounded corners=5pt, fill=offchain!45, inner sep=8pt, fit=(ipfs)(keep), label={[font=\small\bfseries]below:Off chain (Node.js)}] (offz) {};
  \node[draw=boxline, rounded corners=5pt, fill=onchain!55, inner sep=8pt, fit=(fac)(tok)(vault)(rep), label={[font=\small\bfseries]below:On chain (Ethereum, Hardhat local)}] (onz) {};
\end{scope}
\draw[dashed, very thick, black!55] (-1.65,4.1) -- (-1.65,-2.2) node[below, font=\scriptsize]{trust boundary};
% arrows
\foreach \p in {org,mem,adm} \draw[ar] (\p) -- (\p |- ui.north);
\draw[ar] (1.9,0 |- ui.south) -- node[lb]{signed transactions, view calls} (1.9,0 |- fac.north);
\draw[ar] (-3.3,0 |- ui.south) -- node[lb]{upload, download, re hash terms} (-3.3,0 |- ipfs.north);
\draw[ar] (-6.35,0 |- ui.south) |- node[lb, pos=0.3, rotate=90, yshift=5pt]{HTTP API} (idx.west);
\draw[ar, dashed] (circ.west) -- node[lb]{events (read only)} (idx.east);
\draw[ar] (noti) -- (idx);
\draw[ar] (keep.east) -- node[lb, pos=0.45]{public calls} (-0.95,-1.3) -- (-0.95,1.2) -- (circ.west |- 0,1.2);
\draw[ar] (fac) -- node[lb, right]{clone} (circ);
\draw[ar, dashed] (circ) -- node[lb, below=1pt]{delegates} (impl);
\draw[ar] (circ.north east) -- node[lb, below right, inner sep=0pt]{} (tok.south west);
\draw[ar] (circ.south) ++(-0.7,0) -- (vault.north);
\draw[ar] (circ.south) ++(0.3,0) -- (order.north);
\draw[ar] (circ.south) ++(1.2,0) -- (rep.north);
\end{tikzpicture}
\caption{Architecture: seven contracts, four off chain services, one browser app (blue: on chain, amber: off chain)}
\label{fig:architecture}
\end{figure}

\subsection{Contracts and inheritance}
\code{Circle} inherits the shared lifecycle base, which in turn inherits OpenZeppelin's
\code{Initializable}, \code{AccessControl}, \code{Pausable} and \code{ReentrancyGuard}
\cite{openzeppelin}. The shared modules are composed, not inherited: the factory registers every new
circle with the vault, the draw and the registry, and those modules accept writes only from addresses
holding \code{CIRCLE\_ROLE} (Figure~\ref{fig:contracts}).

\begin{figure}[H]
\centering
\begin{tikzpicture}[
  cl/.style={draw=boxline, align=center, font=\small, minimum height=0.85cm, minimum width=2.75cm, fill=white, inner sep=3pt},
  oz/.style={cl, fill=ozgray, font=\footnotesize},
  inh/.style={-{Triangle[open,length=2.6mm,width=2.4mm]}, thick},
  use/.style={-{Stealth[length=2mm]}, draw=black!70},
  lb/.style={font=\scriptsize, align=center, fill=white, inner sep=1pt}]
\node[oz] (ini) at (-4.8,4.2) {Initializable};
\node[oz] (ac) at (-1.6,4.2) {AccessControl};
\node[oz] (pa) at (1.6,4.2) {Pausable};
\node[oz] (rg) at (4.8,4.2) {ReentrancyGuard};
\node[cl, minimum width=12.6cm, fill=tintviolet] (base) at (0,2.6) {\textit{BaseCircleLifecycle} (abstract): state machine, roles, bounded pause, reentrancy guard, initializer};
\node[cl, minimum width=4.6cm, line width=1pt, fill=tintblue] (circle) at (0.6,1.0) {\textbf{Circle}\\[-1pt]{\scriptsize join, contribute, settle, claim, cancel}};
\node[cl, fill=tintamber] (fac) at (-4.9,1.0) {\textbf{CircleFactory}\\[-1pt]{\scriptsize EIP 1167 clones}};
\node[cl, fill=tintteal] (vault) at (-4.4,-1.2) {\textbf{CollateralVault}\\[-1pt]{\scriptsize deposit, release, slash}};
\node[cl, fill=tintteal] (order) at (-1.0,-1.2) {\textbf{PayoutOrder}\\[-1pt]{\scriptsize commit, reveal, finalize}};
\node[cl, fill=tintteal] (rep) at (2.4,-1.2) {\textbf{ReputationRegistry}\\[-1pt]{\scriptsize joins, completions, defaults}};
\node[cl, fill=tintgreen] (tok) at (5.6,-1.2) {\textbf{MockStablecoin}\\[-1pt]{\scriptsize ERC20, 6 decimals}};
\foreach \x in {ini,ac,pa,rg} \draw[inh] (\x.south) -- (\x.south |- base.north);
\draw[inh] (circle.north) -- (circle.north |- base.south);
\draw[use] (fac) -- node[lb, above, yshift=2pt]{clones} (circle);
\draw[use] (circle.south) ++(-1.8,0) -- node[lb, pos=0.5]{CIRCLE\_ROLE} (vault.north);
\draw[use] (circle.south) ++(-0.8,0) -- (order.north);
\draw[use] (circle.south) ++(0.9,0) -- (rep.north);
\draw[use] (circle.south) ++(1.9,0) -- node[lb, pos=0.5]{transfers} (tok.north);
\draw[use, dashed] (fac.south) -- node[lb, left]{registers} (vault.north west);
\end{tikzpicture}
\caption{Inheritance (open triangles) and calls between contracts (arrows)}
\label{fig:contracts}
\end{figure}

\begin{table}[H]
\centering\small
\caption{The seven contracts}
\label{tab:contracts}
\begin{tabularx}{\textwidth}{@{}P{3.6cm}P{2cm}L@{}}
\toprule
Contract & Kind & Responsibility \\
\midrule
\code{BaseCircleLifecycle} & abstract & Lifecycle enum and transition guards; organizer, member and pauser roles; bounded pause whose time extends the round deadline; reentrancy guard; initializer instead of constructor. \\
\code{Circle} & clone & One committee: join, contribute, settle each round, pay out, handle defaults, complete or cancel. \\
\code{CircleFactory} & contract & Deploys one \code{Circle} implementation and creates every circle as a 45 byte EIP 1167 clone, initialized and registered in the same transaction. \\
\code{CollateralVault} & contract & Holds every deposit, keyed by (circle, member); computes the position based amount; deposit, release and slash for registered circles only. \\
\code{PayoutOrder} & contract & Commit and reveal draw; non revealers are left out; the shuffle seed is fixed by the block holding the last reveal. \\
\code{ReputationRegistry} & contract & Counts joins, completions and defaults per address across circles; a circle can refuse members with past defaults. History is emitted as events. \\
\code{MockStablecoin} & contract & Test ERC20 with 6 decimals and a public faucet. \\
\bottomrule
\end{tabularx}
\end{table}

\subsection{On chain versus off chain data}
\label{sec:datamap}
The rule is that contract correctness never depends on an off chain service. Table~\ref{tab:datamap}
lists where each piece of data lives and why.

\begin{table}[H]
\centering\small
\caption{Data map}
\label{tab:datamap}
\begin{tabularx}{\textwidth}{@{}P{3.4cm}P{3.4cm}L@{}}
\toprule
Data & Where & Why, and how integrity is kept \\
\midrule
Deposits & On chain: \code{CollateralVault} & Money is held by code; only registered circles can move it. \\
Contributions, pots, credits & On chain: \code{Circle} & Pull payments behind a reentrancy guard. \\
Rules (amount, size, durations, join rule) & On chain: \code{Circle} configuration & Set once in \code{initialize}; cannot change after people join. \\
State, round, who paid & On chain & Decides who may do what. \\
Commitments, reveals, order & On chain: \code{PayoutOrder} & Fairness must be checkable by every member. \\
Reputation counters & On chain: \code{ReputationRegistry} & Shared across circles; only registered circles write. \\
Terms document & Off chain: IPFS store & Large and human readable; only its CID is on chain and the browser re hashes the download. \\
Display names & Off chain: indexer database & Personal data never goes on chain; each name is signed by the member's wallet. \\
History, lists, reminders & Off chain: indexer and notifier & Convenience; rebuilt from events at any time; the UI falls back to the chain. \\
Draw secret before reveal & Member's browser & Must stay secret; checked against the on chain commitment before reveal. \\
Private keys & Never stored by us & Wallet only; no keys in the repository. \\
\bottomrule
\end{tabularx}
\end{table}
